\section{Related Work}
\label{sec:related-work}

The diversity of existing pub/sub approaches illustrates that there is no one-size-fits-all solution when it comes to designing and implementing a pub/sub system \cite{FacesOfPubSub, XlPubSub}.
Existing approaches differ in their underlying architecture, the mechanisms employed for event dissemination, and the way they address scalability, reliability, and fault tolerance. 

\textbf{Application-Level Multicast (ALM)} approaches primarily focus on creating a separate, highly efficient network overlay for each topic. Scribe~\cite{Scribe} and Bayeux~\cite{Bayeux} construct multicast trees on top of Distributed Hash Tables (DHTs), while Zigzag~\cite{ZigZag} and NICE~\cite{Nice} employ hierarchical clustering. 
Although these approaches facilitate efficient message dissemination, they introduce additional overhead for the construction and maintenance of multicast trees.
Furthermore, tree-based structures introduce potential single points of failure, while message dissemination may involve nodes that are not subscribers of the corresponding topic.

\textbf{Distributed pub/sub} approaches prioritize scalability as the number of nodes and topics grows. 
Gossipsub~\cite{GossipSub} maintains a mesh network for each topic, combining push-based dissemination with gossip-based recovery of missed messages.
It additionally employs a peer-scoring mechanism to evaluate neighboring nodes and improve resilience against malicious behavior.

Poldercast~\cite{Poldercast} combines topic-specific ring topologies with additional random links to achieve efficient and reliable dissemination. 
Its architecture employs Cyclon for random peer sampling and Vicinity for selecting neighbors based on common subscriptions, allowing connections to be reused across multiple topics. 
Our approach follows a similar layered architecture, but introduces a topic-oriented navigation mechanism to facilitate subscriber discovery across topics. 
Furthermore, our framework incorporates secure peer sampling, blockchain-based topic administration, and a separate persistence layer for reliable event retrieval.

Spidercast~\cite{Spidercast} similarly focuses on constructing topic-connected overlays, employing two neighbor-selection strategies: favoring peers with the largest number of common subscriptions, and maintaining random links to peers sharing at least one topic. 
This combination aims to preserve topic connectivity and small dissemination paths while keeping the number of maintained connections moderate.

Tera~\cite{Tera}, on the other hand, focuses on locating subscribers of specific topics rather than optimizing dissemination within a topic. 
It employs a distributed lookup mechanism based on random peer sampling and random walks, allowing nodes to discover access points for topics without prior knowledge of their participants.

\textbf{Centrally managed pub/sub} approaches rely on dedicated servers or brokers for message dissemination. 
Dynamoth~\cite{Dynamoth} employs dynamic resource allocation and load-balancing strategies to accommodate changes in communication workloads, while Dynatops~\cite{Dynatops} utilizes a broker-based architecture and centralized reconfiguration to minimize dissemination costs. 
Unlike distributed approaches, these approaches depend on designated infrastructure, requiring additional server resources as the communication workload increases.